\subsection{The proof strategy}
The ten supporting hyperplanes may move in all four coordinate directions,
so the nearby bodies include polytopes that are no longer products.
Write
\[
 K(a)=\{x\in\mathbb R^4:a_i\cdot x\leq1,\ i=1,\ldots,10\},
 \qquad a=(a_1,\ldots,a_{10}),
\]
and let $a_0$ be the row list of $K_0$.

We construct twenty-six smooth upper bounds $U_j$ for the systolic ratio,
all equal to it at $a_0$. If one of these bounds decreases, the ratio
decreases with it. Different bounds handle different directions.
The proof first identifies their derivatives and the symmetry directions.
A linear-algebra lemma then gives a uniform negative slope transverse to
symmetry. Only at the end do we pass from these derivatives to a finite
neighborhood and recover the equality cases.

\subsection{Local coordinates and volume}
\begin{lemma}[A simple row neighborhood]\label{lem:hko-row-neighborhood}
There is a neighborhood $\mathcal N$ of $a_0$ in $(\mathbb R^4)^{10}$
on which $K(a)$ is bounded and simple, with the same vertex--facet
incidences as $K_0$. Every sufficiently Hausdorff-close ten-facet polytope
has a labelled presentation with $a\in\mathcal N$, and its rows tend to
$a_0$ as the body tends to $K_0$.
\end{lemma}
\begin{proof}
The origin stays in the interior under sufficiently small Hausdorff
perturbations of $K_0=K(a_0)$. One way to see this is through support
functions: $h_K(u)=\max_{x\in K}u\cdot x$. Write $B$ for the Euclidean
unit ball. If $rB\subset K_0$ and the Hausdorff distance from $K$ to $K_0$
is less than $\epsilon<r$, then
$h_K(u)\geq h_{K_0}(u)-\epsilon\geq r-\epsilon$ for every unit vector
$u$. Thus $(r-\epsilon)B\subset K$. On such a neighborhood polarity is
continuous: the radial function, the distance from zero to the boundary
along a unit direction, of $K^\circ$ is $1/h_K$ on the unit
sphere, and the common positive lower bound on $h_K$ makes these reciprocals
converge uniformly. The polars are uniformly bounded convex bodies, so this
also gives Hausdorff convergence.

The polar $K_0^\circ$ has the ten vertices $a_{0,i}$. Choose disjoint small
balls about them. Each vertex is uniquely exposed by some linear functional
$\ell_i$. For polars sufficiently close to $K_0^\circ$, every maximizer
of $\ell_i$ lies in the corresponding ball. Otherwise a convergent sequence
of maximizers outside that ball would limit to a maximizer on $K_0^\circ$
different from $a_{0,i}$. A maximizing face contains a vertex. A polytope
$K$ with exactly ten facets has a polar with exactly ten vertices; hence
there is exactly one polar vertex in each ball and none elsewhere. Labelling
these vertices $a_i$ gives
\[
 K=\{x:a_i\cdot x\leq1,\ i=1,\ldots,10\},\qquad a_i\longrightarrow a_{0,i}
 \quad\text{as }K\longrightarrow K_0.
\]

For completeness, the face structure is stable in this neighbourhood.
The base body is a product of two polygons and is simple: precisely four
facets meet at each vertex, with independent normals. Their four equations
therefore determine a smoothly varying intersection under a small change
of $a$. All the remaining inequalities have strict slack at that
intersection, so it remains a vertex with exactly those incident facets.

No new vertex can appear. The convex hull of the base normals contains a
ball about zero; the convex hull of sufficiently close normals contains a
smaller common ball. Consequently the bodies $K(a)$ are uniformly bounded.
If arbitrarily close parameters had a new vertex, choose four independent
active constraints there. There are finitely many four-element subsets,
so along a subsequence the same subset is active and the vertices converge
to a point of $K_0$ where those four constraints are equalities. Simplicity
forces them to be the incident facets of one of the original vertices.
Their intersection is the continuation already constructed, contradicting
newness. This argument also excludes a new vertex arising from a subset
whose normals are dependent at the base point; no invertibility assumption
for such a subset is needed.

The vertex--facet incidences therefore remain unchanged. In particular
each inequality still defines a facet. Conversely,
small changes of the ten normals give Hausdorff-small changes of $K(a)$:
their convex hulls converge, contain a common ball about zero, and polarity
is continuous there.

\end{proof}

\begin{lemma}[Smooth volume and its derivative]\label{lem:hko-smooth-volume}
On a sufficiently small neighborhood $\mathcal N$ as above, $V(a)=
\operatorname{vol}_4K(a)$ is smooth and positive. If $F_i(a)$ has centroid
$\bar x_i(a)$, then
\[
 DV(a)[h]=-\sum_{i=1}^{10}
 \frac{\operatorname{vol}_3(F_i(a))}{\|a_i\|}
 \langle h_i,\bar x_i(a)\rangle.
\]
\end{lemma}
\begin{proof}
The fixed incidences express each vertex smoothly in the rows. A fixed
boundary triangulation then expresses volume smoothly in those vertices.
On the moving facet, differentiating
$\langle a_i+\epsilon h_i,x(\epsilon)\rangle=1$ gives outward normal
velocity $-\langle h_i,x\rangle/\|a_i\|$. Integrating this velocity over
the facets and using their centroids gives the formula.
\end{proof}

\subsection*{The concrete base point and its volume}
Put $t=\tan(\pi/5)$, $s=(5-t^2)/2=\sqrt5$, and
\[
 t^4-10t^2+5=0,\quad 0<t<1,\qquad
 \alpha=(3-s)/2,\quad b=t(1+s)/2,\quad d=s-1.
\]
The selected root is unique in $(0,1)$: the polynomial is strictly decreasing
there and changes sign. For the circumradius-one pentagon with vertex $(1,0)$,
the edge normals have arguments $(2k+1)\pi/5$ and lengths $\sec(\pi/5)=d$.
Rotating the second factor by $-\pi/2$ gives the following ten rows $a_{0,i}$, in
$(q_1,q_2,p_1,p_2)$ order:
\[
\begin{pmatrix}
1&t&0&0\\ -\alpha&b&0&0\\-d&0&0&0\\-\alpha&-b&0&0\\1&-t&0&0\\
0&0&t&-1\\0&0&b&\alpha\\0&0&0&d\\0&0&-b&\alpha\\0&0&-t&-1
\end{pmatrix}.
\]
These rows give the normalized facet presentation of $K_0$.
Writing $\theta=\pi/5$, the area of one pentagon is the sum of five
triangles with central angle $2\theta$, hence
$B=\frac52\sin(2\theta)=5\sin\theta\cos\theta$.
Fubini therefore gives
\[
 V_0=B^2=25\sin^2\theta\cos^2\theta
 =\frac{25(5+s)}{32}=\frac{25(15-t^2)}{64}.
\]
To specialize Lemma~\ref{lem:hko-smooth-volume}, observe that $F_i$ is an edge times the other centered pentagon. Its centroid is
$\bar x_i=a_{0,i}/\|a_{0,i}\|^2$, its three-volume is $2\sin\theta\,B$,
and $\|a_{0,i}\|=\sec\theta$. Consequently
\[
 DV_0[h]=-\mu\sum_{i=1}^{10}a_{0,i}\cdot h_i,\qquad
 \mu=10\sin^2\theta\cos^4\theta=\frac{25}{32}+\frac{5s}{16}.
\]
This gives the volume derivative for every row perturbation at the base point.
